\documentclass[10pt]{article}
\usepackage[margin=25mm]{geometry}
\usepackage{fontspec}
\usepackage{unicode-math}
\usepackage{amsmath,booktabs,longtable,array,graphicx,xurl}
\usepackage{calc}
\usepackage{microtype}
\usepackage{fancyvrb}
\usepackage{hyperref}
\providecommand{\tightlist}{\setlength{\itemsep}{0pt}\setlength{\parskip}{0pt}}
\providecommand{\real}[1]{#1}
\hypersetup{colorlinks=true,linkcolor=blue,citecolor=blue,urlcolor=blue}
\setlength{\parskip}{0.45em}
\setlength{\parindent}{0pt}
\setlength{\emergencystretch}{3em}
\newcommand{\doctype}{Technical report}
\newcommand{\docversion}{v0.10.006}
\newcommand{\docdate}{30 October 2026}
\newcommand{\versiondoi}{10.5281/zenodo.23826265}
\newcommand{\conceptdoi}{10.5281/zenodo.23783845}
\newcommand{\docrepo}{https://github.com/fsantibanezleal/CAOS\_OreFlow}
\newcommand{\headerblock}{%
\begin{center}
\fbox{\parbox{0.94\textwidth}{\small\raggedright
\textbf{\doctype} \quad$\cdot$\quad \docversion, \docdate \quad$\cdot$\quad Licence CC BY 4.0\\[2pt]
CAOS open-research programme \quad$\cdot$\quad mineral-processing simulator\\[2pt]
Version DOI \href{https://doi.org/\versiondoi}{\texttt{\versiondoi}} \quad$\cdot$\quad
concept DOI (always latest) \href{https://doi.org/\conceptdoi}{\texttt{\conceptdoi}}\\[2pt]
Not peer reviewed. Source code, data and computational records: \href{https://github.com/fsantibanezleal/CAOS_OreFlow}{\texttt{github.com/fsantibanezleal/CAOS\_OreFlow}}
}}
\end{center}
\vspace{4pt}}
\title{How far does a learned surrogate of a flowsheet simulator transfer across ores? A closed-circuit, size-by-mineral engine, with cases and ore groups held out}
\author{Felipe Santibáñez-Leal, ORCID 0000-0002-0150-3246}
\date{}
\begin{document}
\maketitle
\headerblock
\section*{Abstract}

Learned surrogates of process simulators are often evaluated on random
states from plants they have already seen. We compare that interpolation
setting with leave-one-case-out and leave-one-ore-group-out evaluation
for five models trained on a steady-state mineral-processing simulator.
The simulator tracks mineral mass over 63 size classes, water, closed
grinding circuits and unit-level balance closure; its limitations
against published examples are reported alongside the comparisons. On a
3,072-state design over twelve authored cases, the multilayer perceptron
gives the lowest interpolation recovery RMSE (3.38 percentage points),
but has a 26.24-point mean RMSE when each case is held out. The Gaussian
process has the lowest mean held-out-case RMSE (10.95 points), while the
perceptron has the highest median case-fold R² (0.547). Under ore-group
holdout, the perceptron's median recovery R² is 0.437 across seven
groups, but varies from 0.795 for the five chalcopyrite cases to -10.2
for the two gold-in-pyrite cases. Log-upgrade transfer is poor for every
model. Thus interpolation, case holdout and ore-group holdout answer
different questions, and their results depend on both target and ore
family. The records and tests bind every reported score to the bake. The
cases are authored; these results do not establish performance on
operating plants.

\section{Introduction}\label{introduction}

A surrogate of a flowsheet simulator is attractive for the same reason
as any surrogate: the simulator is slow and the surrogate is fast, so it
can sit inside an optimizer, a sweep or an interactive tool. Whether it
can stand in for the simulator depends on the question it is asked.
Inside the envelope it was trained on, an interpolation score answers
that question. A digital twin is asked something harder: what happens
with a different ore, a harder or finer feed, another plant. That is
transfer, and an interpolation score says nothing about it.

This work measures both on the same models, with a simulator whose
answers are known exactly and can be reproduced. Its contributions are
three:

\begin{enumerate}
\def\labelenumi{\arabic{enumi}.}
\tightlist
\item
  an open flowsheet engine that represents streams by size and mineral,
  solves closed circuits, audits every balance from its output streams,
  is checked against published examples with every miss stated, and is
  reproduced by a line-by-line second implementation;
\item
  a test protocol that scores learned surrogates of it by interpolation,
  by leaving one case out and by leaving one ore group out over twelve
  authored plants in four circuit families, with interval coverage and
  an out-of-envelope guard;
\item
  the findings that holding out one plant is not a transfer test when
  other plants share its ore, that transfer failure is concentrated in
  the plants unlike the others, that the ranking of models changes
  between the protocols and with the statistic that summarizes them, and
  that the guard misses shifts of inputs its training features do not
  correlate with.
\end{enumerate}

\section{The engine}\label{the-engine}

The record reproduces all 96 committed states within 1e-6 in the browser
port and all 96 committed states close within 1e-9 at every unit in the
Python engine. The port shares constants and is a translation check, not
an independent validation of physical assumptions.

\subsection{Representation}\label{representation}

Every stream holds, for each mineral, a vector of mass flows (t/h) over
one grid of 63 size classes on a fourth-root-of-two progression from 150
mm to a pan below 3.24 µm, and a water flow. Assays are computed from
mineral masses and the minerals' element contents (from their formulas
and the IUPAC atomic weights); no grade is stored, so a grade changes
only when a mass does. Liberation is described per valuable mineral by
size (King 1979): a liberated fraction, and composites whose content and
host are declared.

\subsection{Comminution}\label{comminution}

The crusher is the Whiten form (Duarte et al.~2021), with a strictly
lower-triangular breakage matrix \texttt{B} and a classification
function \texttt{C} set by the closed-side setting:

\begin{verbatim}
p = (I - C)(I - B C)^-1 f
\end{verbatim}

The ball mill is an energy-specific population balance (Herbst and
Fuerstenau 1980) in the form of the Moly-Cop tools: three perfect mixers
of volume fractions 0.70, 0.15 and 0.15 share the breakage operator
\texttt{D\ =\ (I\ -\ B)\ diag(S\^{}E)}, so the mill's inverse transfer
at an energy per pass \texttt{e} is a cubic polynomial in \texttt{D},
and the closed circuit with the cyclone returning a fraction
\texttt{C\_i} of each class solves

\begin{verbatim}
(T^-1(e) - diag(C)) p = f,        T^-1(e) = I + c1 e D + c2 e^2 D^2 + c3 e^3 D^3
\end{verbatim}

per mineral. Two scalar unknowns, the energy per pass that meets the
target P80 and the host cut that holds the design circulating load, are
found by the Illinois method. The energy per pass is per tonne of mill
feed, so the mill draws \texttt{e\ m}, with \texttt{m} the fresh feed
plus the cyclone underflow (less what a gravity bleed takes from it);
when that exceeds the installed power, the circuit runs at installed
power, \texttt{e\ m(e)\ =\ P\_inst} solved as a fixed point, with a
coarser product, and says so. The reported grinding energy is that net
energy per tonne of fresh feed, \texttt{e\ m\ /\ F}; the crushing energy
is Bond's with the crushing work index, and the Bond requirement
\texttt{E\ =\ 10\ W\_i\ (P\^{}-1/2\ -\ F\^{}-1/2)} gives the operating
work index and the efficiency ratio beside it, with Rittinger and Kick
calibrated to agree with Bond at a reference reduction and never added
to it (Bond 1952; GMG 2021).

\subsection{Classification and
separation}\label{classification-and-separation}

Hydrocyclones follow Plitt (1976): a Rosin-Rammler partition
\texttt{y(d)\ =\ R\_f\ +\ (1\ -\ R\_f)(1\ -\ exp(-ln\ 2\ (d/d50c)\^{}m))}
with a water bypass \texttt{R\_f} from the water balance and a cut
corrected for each mineral's density, so dense minerals classify as
coarser particles and circulate. Plitt's sizing of the cyclone cluster
is uncalibrated and fails its published check (section 3), so the
cluster's count and pressure are reported as an estimate, never as a
result. Flotation banks are perfect mixers in series; the rate constant
of each particle class follows from the bubble surface area flux
\texttt{S\_b\ =\ 6\ J\_g\ /\ D\_32} (Gorain et al.~1997, 1999) and a
size window (Trahar 1981), a composite floats at the liberated rate
scaled by its exposed surface, \texttt{P\_comp\ =\ P\_V\ phi\^{}(2/3)}
with \texttt{phi} the valuable mineral's volume share, entrainment
follows Savassi et al.~(1998), and cleaner and recleaner tails return to
the preceding bank. Each cell's residence is its pulp volume, less the
gas hold-up, over the flow that leaves it as tail, which depends on what
the cell floats, so it is solved per cell as a fixed point:

\begin{verbatim}
r_j = (k tau_j + ENT w_j) / (1 + k tau_j + ENT w_j)
tau_j = V (1 - eps_g) / Q_tail,j
R_N = 1 - prod_j (1 - r_j)
\end{verbatim}

Gold follows the gravity-recoverable gold (GRG) model of Laplante,
Woodcock and Noaparast as Vincent (1997) states it: the gold is a GRG
component with its own size distribution, from a GRG test, plus a
non-GRG component; the GRG's cyclone cut takes a density correction with
an exponent of 1.0, fitted to measured GRG partitions where Stokes would
give 0.5; the gravity unit on a bleed of the cyclone underflow, or of
the mill discharge, recovers GRG by size at a per-pass recovery
\texttt{R(d)\ =\ R\_max\ (1\ -\ exp(-(d/x\_g)\^{}2))}; and grinding
moves GRG to non-GRG at a slowed breakage rate (Banisi et al.).
Low-intensity magnetic drums capture liberated magnetite and composites
by their magnetite content; desliming cyclones discard slimes ahead of
flotation, with the water split following the solids sent down at a
declared underflow density.

\subsection{Balance closure and the
contract}\label{balance-closure-and-the-contract}

An independent recomputation calculates, from the output streams alone,
the closure of every mineral, every species and the water at every unit
and for the circuit, and of every size class at the units that do not
break particles. The operating point is thirteen inputs declared once,
with units, per-case bounds, steps, the families they apply to, a
cross-field rule and bilingual messages; the exported declaration is
interpreted identically by the engine's validator, the web service and
the browser, so a state is accepted or rejected with the same code
everywhere. A state the engine cannot solve, such as a classifier cut at
which the mill cannot draw its installed power, is refused with a reason
rather than answered.

\section{Verification}\label{verification}

\textbf{Balances.} All 96 committed states close within 1e-9 relative at
every unit and for the circuit, as recomputed from the stored streams by
the artifact checks.

\textbf{The second implementation.} A line-by-line TypeScript port runs
the engine in the browser. Re-simulating every committed state from its
own definition, it matches every metric, stream record, curve and
kinetic record within 1e-6 relative, with the same flags. It is a
translation by the same author that reads the same constant,
atomic-weight and mineral files, so it catches translation errors and
not modelling errors. An independent verification of the 0.08.001
implementation against its equations, one reviewer per part of the
engine and every finding verified, found defects in both engines, among
them the composite rate, the cell residence, the mill power basis and
the desliming water split; each was fixed, the engine's in both
languages, with a test that fails on the earlier code, and every number
below is from the fixed engine.

\textbf{Published examples.} Each is labelled as a published example,
not as plant data:

{\def\LTcaptype{none} % do not increment counter
\small
\begin{longtable}[]{@{}
  >{\raggedright\arraybackslash}p{(\linewidth - 4\tabcolsep) * \real{0.3333}}
  >{\raggedright\arraybackslash}p{(\linewidth - 4\tabcolsep) * \real{0.3333}}
  >{\raggedright\arraybackslash}p{(\linewidth - 4\tabcolsep) * \real{0.3333}}@{}}
\toprule\noalign{}
\begin{minipage}[b]{\linewidth}\raggedright
Example
\end{minipage} & \begin{minipage}[b]{\linewidth}\raggedright
Published
\end{minipage} & \begin{minipage}[b]{\linewidth}\raggedright
Engine
\end{minipage} \\
\midrule\noalign{}
\endhead
\bottomrule\noalign{}
\endlastfoot
Moly-Cop BallSim base case (504 t/h, every published input) & net 7.71
kWh/t, gross 8.56; cut 183.3 µm; 6 cyclones at 53 kPa & net 7.30 kWh/t,
5.2\% below, inside the 20\% tolerance set before the first run; cut 188
µm; Plitt's uncalibrated sizing asks for 2 cyclones at 816 kPa \\
GMG worked examples of the operating work index & 14.4 and 11.7 kWh/t &
14.38 and 11.71 kWh/t \\
Laplante and Staunton gravity example, run like for like (80.4\% GRG,
the unit on 10 to 60\% of the mill discharge) & GRG recovery 79.8 to
95.9\%, GRG circulating load 2016 to 413\% & GRG recovery 69.5 to
90.7\%, GRG circulating load 696 to 84\%: outside the 5-point and 35\%
tolerances \\
Zandrivierspoort magnetite (a different ore; the 45 µm product is a
regrind) & 64.9 and 69.0\% Fe at 75 and 45 µm & 63.8 and 67.8\% Fe: a
step of 4.0 points against 4.1, at levels 1.1 and 1.2 points below \\
\end{longtable}
}

The gravity example misses because the measured GRG of the case's ore
(Snip) is very fine: 9 to 31\% of the GRG, almost all finer than about
37 µm, leaves by the overflow; without the GRG finer than 25 µm the GRG
recovery comes within 1.1 points from the 20\% row on, and the
circulating load stays 50 to 69\% low. The magnetite example checks a
direction and a proportion, not a level, because its ore is not the
engine's case.

\section{The test}\label{the-test}

\subsection{The plants}\label{the-plants}

Twelve authored cases in four circuit families: nine copper, nickel,
zinc and refractory-gold rougher-cleaner flotation circuits
(\texttt{rougher}), free-milling gold with a gravity bleed
(\texttt{gravity\_rougher}), fine magnetite with magnetic drums instead
of flotation (\texttt{magnetic}), and phosphate with desliming ahead of
flotation (\texttt{deslime\_rougher}). Every parameter carries its unit
and either a source or the label authored; none is fitted to a plant.
Five of the twelve are copper sulphide plants on chalcopyrite, and the
two gold plants carry most of their gold in pyrite (55\%, beside
electrum in the free-milling ore and arsenopyrite in the refractory
one); the other five each have a payable and dominant carrier no other
case shares.

\subsection{The design, the features and the
targets}\label{the-design-the-features-and-the-targets}

For each case a scrambled Sobol sequence draws 256 states the contract
accepts over the case's whole envelope and two ore factors (liberation
size and floatability within their uncertainty half-widths), 3072 in
all; the engine simulates every state. Each state is described by 22
physical features computed before simulation: the payable fraction,
throughput per installed megawatt, the grind, circulating load, water,
crusher setting and work index, the payable carriers' share-weighted
liberation size, composite content, density, floatability and dose
ratio, the rougher gas velocity, cells and volume per t/h, the gravity
bleed, the desliming cut and four circuit descriptors. The case identity
is never a feature. The targets are recovery (\%), the log10 upgrade
ratio (concentrate over head grade) and the total specific energy
(kWh/t).

\subsection{The models and the
protocols}\label{the-models-and-the-protocols}

Ridge (alpha 1); a random forest (300 trees; Breiman 2001); histogram
gradient boosting (300 iterations; Friedman 2001); a Gaussian process
with a squared-exponential kernel with one length scale per feature and
a noise term, on a seeded subsample of 500 training rows (Rasmussen and
Williams 2006); and a perceptron of two hidden layers of 64 units,
trained by full-batch AdamW with early stopping on a 15\% validation
split (Loshchilov and Hutter 2019). Features and targets are
standardized on the training rows of each split. Because the Gaussian
process sees fewer rows, the other models are also refitted on the same
500 rows; and because one training seed is one draw, the perceptron is
retrained with four more seeds on every split.

\emph{Interpolation} holds out 20\% of the states inside every case
(2460 training, 612 test rows). \emph{Leave one case out} holds out each
case in turn; the models that predict a case never saw any of its
states, but they may have seen plants on the same ore. \emph{Leave one
ore group out} holds out every plant that shares the held-out plant's
payable and dominant carrier mineral: the five chalcopyrite plants
together (1792 training, 1280 test rows), the two gold plants together
(2560 training, 512 test rows), and each of the other five alone, where
it is the same fold as leaving one case out; seven groups in all. An
autoencoder (16, 6 and 16 units) trained on the features gives a guard:
its threshold is the 99th percentile of the reconstruction errors of its
validation rows; its false-alarm rate is measured on the 612 held-out
in-envelope states, and its false-accept rate on 11 016 probes, each a
held-out state pushed past the training range of one continuous feature,
at six distances from 0.02 to 1.0 of the range, above the maximum and
below the minimum.

\section{Results}\label{results}

\subsection{Interpolation and held-out cases answer different
questions}\label{interpolation-and-held-out-cases-answer-different-questions}

\textbf{Table 1.} Recovery, RMSE in percentage points. The last column
is the median R² over seven held-out ore groups.

{\def\LTcaptype{none} % do not increment counter
\small
\begin{longtable}[]{@{}
  >{\raggedright\arraybackslash}p{(\linewidth - 12\tabcolsep) * \real{0.1111}}
  >{\raggedleft\arraybackslash}p{(\linewidth - 12\tabcolsep) * \real{0.1481}}
  >{\raggedleft\arraybackslash}p{(\linewidth - 12\tabcolsep) * \real{0.1481}}
  >{\raggedleft\arraybackslash}p{(\linewidth - 12\tabcolsep) * \real{0.1481}}
  >{\raggedleft\arraybackslash}p{(\linewidth - 12\tabcolsep) * \real{0.1481}}
  >{\raggedleft\arraybackslash}p{(\linewidth - 12\tabcolsep) * \real{0.1481}}
  >{\raggedleft\arraybackslash}p{(\linewidth - 12\tabcolsep) * \real{0.1481}}@{}}
\toprule\noalign{}
\begin{minipage}[b]{\linewidth}\raggedright
Model
\end{minipage} & \begin{minipage}[b]{\linewidth}\raggedleft
Interpolation RMSE
\end{minipage} & \begin{minipage}[b]{\linewidth}\raggedleft
Interpolation R²
\end{minipage} & \begin{minipage}[b]{\linewidth}\raggedleft
Leave one case out: mean RMSE
\end{minipage} & \begin{minipage}[b]{\linewidth}\raggedleft
worst case
\end{minipage} & \begin{minipage}[b]{\linewidth}\raggedleft
median R²
\end{minipage} & \begin{minipage}[b]{\linewidth}\raggedleft
Leave one ore group out: median R²
\end{minipage} \\
\midrule\noalign{}
\endhead
\bottomrule\noalign{}
\endlastfoot
Ridge & 9.46 & 0.617 & 13.68 & 23.05 & -0.127 & -0.774 \\
Random forest & 6.49 & 0.820 & 13.16 & 48.86 & 0.359 & 0.129 \\
Histogram gradient boosting & 4.60 & 0.909 & 12.66 & 56.61 & 0.512 &
0.287 \\
Gaussian process & 5.34 & 0.878 & 10.95 & 27.91 & 0.481 & 0.171 \\
Perceptron & 3.38 & 0.951 & 26.24 & 239.51 & 0.547 & 0.437 \\
\end{longtable}
}

The perceptron has the lowest interpolation RMSE (3.38 points) and
highest pooled interpolation R² (0.951), while its mean
leave-one-case-out RMSE is highest (26.24 points) because the magnetite
fold is difficult. The Gaussian process has the lowest mean
leave-one-case-out error (10.95 points); the perceptron has the highest
median fold R² (0.547). The five-seed record gives the MLP interpolation
RMSE range of 3.00 to 3.54 points and mean held-out-case RMSE of 25.05
to 30.10 points. These are different summaries and support no single
model ranking for every use.

\subsection{Case-level errors identify where interpolation
fails}\label{case-level-errors-identify-where-interpolation-fails}

\textbf{Table 2.} Recovery RMSE (points) on each held-out case.

{\def\LTcaptype{none} % do not increment counter
\small
\begin{longtable}[]{@{}
  >{\raggedright\arraybackslash}p{(\linewidth - 10\tabcolsep) * \real{0.1304}}
  >{\raggedleft\arraybackslash}p{(\linewidth - 10\tabcolsep) * \real{0.1739}}
  >{\raggedleft\arraybackslash}p{(\linewidth - 10\tabcolsep) * \real{0.1739}}
  >{\raggedleft\arraybackslash}p{(\linewidth - 10\tabcolsep) * \real{0.1739}}
  >{\raggedleft\arraybackslash}p{(\linewidth - 10\tabcolsep) * \real{0.1739}}
  >{\raggedleft\arraybackslash}p{(\linewidth - 10\tabcolsep) * \real{0.1739}}@{}}
\toprule\noalign{}
\begin{minipage}[b]{\linewidth}\raggedright
Held-out case
\end{minipage} & \begin{minipage}[b]{\linewidth}\raggedleft
Ridge
\end{minipage} & \begin{minipage}[b]{\linewidth}\raggedleft
Random forest
\end{minipage} & \begin{minipage}[b]{\linewidth}\raggedleft
Gradient boosting
\end{minipage} & \begin{minipage}[b]{\linewidth}\raggedleft
Gaussian process
\end{minipage} & \begin{minipage}[b]{\linewidth}\raggedleft
Perceptron
\end{minipage} \\
\midrule\noalign{}
\endhead
\bottomrule\noalign{}
\endlastfoot
Soft copper porphyry & 12.54 & 11.59 & 9.47 & 11.09 & 5.98 \\
Hard copper porphyry & 9.62 & 6.72 & 4.82 & 5.07 & 1.72 \\
Free-milling gold with gravity & 8.55 & 12.48 & 10.53 & 15.83 & 12.76 \\
Fine magnetite concentration & 21.01 & 48.86 & 56.61 & 27.91 & 239.51 \\
Nickel sulphide with serpentine slimes & 13.47 & 8.49 & 6.80 & 6.66 &
6.61 \\
Phosphate with clay slimes & 15.85 & 10.04 & 10.05 & 10.83 & 8.93 \\
Copper-molybdenum bulk flotation & 9.38 & 6.12 & 4.03 & 6.33 & 1.91 \\
Oxide copper: assumed sulphidized response & 23.05 & 14.12 & 15.19 &
12.94 & 11.60 \\
Zinc sulphide & 10.87 & 11.88 & 10.18 & 8.49 & 7.26 \\
Copper ore with clay & 10.69 & 5.99 & 6.06 & 9.47 & 6.57 \\
Low-grade copper at high throughput & 9.68 & 6.25 & 4.48 & 6.41 &
2.23 \\
Refractory gold in sulphides & 19.44 & 15.41 & 13.77 & 10.32 & 9.86 \\
\end{longtable}
}

The copper cases are not uniformly easy: holding out the soft porphyry
costs the perceptron 6.0 points, while its error on the hard porphyry,
molybdenum and low-grade copper cases is 1.7 to 2.2 points. Mixed
high-clay copper is 6.6 points. The magnetite fold remains the largest
perceptron error (239.5 points); the tree models also extrapolate poorly
there. These model errors are measured against this authored simulator,
not a plant assay.

\subsection{Whole-ore-group transfer}\label{whole-ore-group-transfer}

When an ore group is held out, median recovery R² across seven groups
ranges from -0.774 for ridge to 0.437 for the perceptron. The MLP scores
0.795 on the five chalcopyrite cases held out together and -10.2 on the
two gold-in-pyrite cases; the group result is strongly dependent on ore
family. Three of five models have a positive median recovery R² across
groups, but no such aggregate removes the poor individual-group results.
The release therefore reports per-case and per-group errors alongside
pooled scores.

\subsection{Energy and upgrade
targets}\label{energy-and-upgrade-targets}

The one-case-out median R² for total specific energy ranges from 0.652
(ridge) to 0.952 (gradient boosting); under ore-group holdout it ranges
from 0.617 (ridge) to 0.907 (Gaussian process). Log-upgrade median R² is
negative for every model under group holdout, so the recovery transfer
result does not generalize to grade-upgrade prediction. The held-out
protocols and target definitions remain separate in the record.

\subsection{Interval coverage}\label{interval-coverage}

The Gaussian process intervals are compared with their empirical
coverage in the committed record. Coverage is target- and
split-specific; the point estimates are not used as a substitute for
uncertainty validation. The authored-plant split does not establish
calibrated intervals on real plants.

\subsection{The guard}\label{the-guard}

The autoencoder guard is evaluated both on in-envelope states and on
deliberately shifted probes. Its per-feature and per-distance acceptance
rates are recorded in \texttt{learning.json}; a low reconstruction error
is not treated as proof that an input shift is physically safe. The
guard is a screening signal, not a process constraint or a substitute
for the operating contract.

\section{The method records}\label{the-method-records}

The engine produces, for each variant, records that do not depend on
learning:

\begin{itemize}
\tightlist
\item
  \textbf{Kinetic lumping.} Five lumped models fitted to a virtual batch
  test of the rougher feed, projected to the bank under the same
  residence distribution and compared with the exact bank recovery by
  true flotation, over 88 flotation states: mean absolute projection
  errors of 1.27 points for the gamma model, 1.81 for Kelsall, 2.80 for
  Klimpel, 4.76 for the stretched exponential and 7.56 for first order
  (Polat and Chander 2000; Vinnett and Waters 2025). The first-order
  projection falls below the exact result at every nominal state: a
  batch curve that is a mixture of rates is not one exponential.
\item
  \textbf{Constrained optimization.} A generalized pattern search with a
  progressive barrier from six starts (Torczon 1997; Audet and Dennis
  2006, 2009) finds a point within the grade, power and water
  constraints for 94 of the 96 variants; the two exceptions are the
  magnetite circuit's harder ore and higher throughput, where the grind
  is the only decision and the mill is already at installed power. The
  learned lane screens the search step (Booker et al.~1999): candidates
  are ranked by the perceptron's prediction and reach the engine only
  where the guard accepts their state and the Gaussian process's 95\%
  half-width on recovery is within 5 points, a novelty measure for the
  perceptron rather than an interval on what it ranks; the poll and the
  optimum stay engine results. The screen did not save engine work: over
  the 72 screened variants it used 26,367 engine evaluations versus
  22,487 without screening, 17.25\% more; it saved calls in none, used
  more in 68 and tied in four. Mean absolute recovery error averaged
  0.66 percentage points across variants; across 5,811 proposals its
  mean was 0.71 points, median 0.48 and maximum 38.16. Both procedures
  had the same terminal outcome in 65 variants: 63 identical feasible
  decisions and two matching infeasible outcomes. Seven feasible pairs
  selected different decisions, with recovered-metal differences no
  greater than 0.002653 t/h (0.017\%). Five favored the unscreened
  decision and two the screened decision.
\item
  \textbf{Sensitivity.} Sobol indices of four uncertain ore properties
  at each nominal state (Saltelli et al.~2010; Herman and Usher 2017):
  floatability drives recovery in ten cases, liberation size drives
  concentrate grade in seven and the head grade in the other five, the
  work index drives grinding energy and the head grade drives recovered
  metal in all twelve.
\end{itemize}

\section{Measured data, which calibrate
nothing}\label{measured-data-which-calibrate-nothing}

The engine calibrates nothing and is calibrated by nothing. Measured
data feed three lanes, and one of them also runs through the engine as
inputs.

\textbf{Particles.} The HZDR constructed-case workbook (Pereira et
al.~2021; RODARE 336, CC BY 4.0) keeps its training sheet (68 008
particles with A/B labels) and its test sheet (29 147 particles with
constructed probabilities and the source authors' predictions, but no
observed labels). A sparse logistic model per case and a shared network
(4, 32, 32, 4) are scored against the constructed probabilities:

{\def\LTcaptype{none} % do not increment counter
\begin{longtable}[]{@{}llll@{}}
\toprule\noalign{}
Case & Published reference & L1 logistic & Network \\
\midrule\noalign{}
\endhead
\bottomrule\noalign{}
\endlastfoot
1 & 0.0257 & 0.0131 & 0.0138 \\
2 & 0.0175 & 0.0157 & 0.0181 \\
3 & 0.0207 & 0.0171 & 0.0143 \\
4 (28 484 complete rows) & 0.0371 & 0.1919 & 0.0226 \\
\end{longtable}
}

The linear model beats the network in cases 1 and 2 and fails only in
case 4, where the network does not.

\textbf{Locked-cycle tests.} 52 GeoMet locked-cycle copper recoveries
from 29 drill holes of one deposit (Hoffimann et al.~2022; archived
dataset 7051975, CC BY 4.0), predicted from five assays by the training
mean, ridge, a random forest and a Gaussian process, under five folds of
whole holes and three spatial zones. RMSE ranges from 5.09 to 5.51
points (whole holes) and 5.15 to 5.69 (zones). A paired bootstrap over
complete holes (2000 resamples) shows how little separates them. On the
five fixed hole folds ridge beats the training mean by 0.42 points (95\%
interval 0.02 to 0.82), but that partition sits at the 98.5th percentile
of 200 random hole partitions; averaged over them the gain is 0.17
points (-0.39 to 0.71) and under leave one hole out 0.22 (-0.36 to
0.79), each interval widened for the six model pairs, and under spatial
zones no model separates either. The source has no grind, reagent or
residence information, so it cannot calibrate the engine's controls.

\textbf{The engine on the same samples.} The 52 samples also run through
the engine as inputs, each on its own assays in the soft copper
porphyry's circuit at that case's operating point; nothing in the engine
is fitted to them. The work index comes from the nearest comminution
sample in the hole (42 samples) or the deposit median (10), and spans
15.2 to 26.1 kWh/t against the case's 11.0. Most of the copper is
bornite and chalcocite, allocated by a sulphur-limited normative
mineralogy, one choice in a family of allocations; the two minerals
float at declared ratios to chalcopyrite (0.8, and chalcocite at 1.5
times bornite, a choice the GeoMet data do not bound). Every sample runs
the mill at installed power, and the engine's recovery falls short of
the locked-cycle test by 22.6 points on average (RMSE 24.1 points; 2.9
to 41.1 points below). That gap is mostly the host circuit's size: the
samples are harder than the circuit's design ore, so at 720 t/h the
product is coarse. With the circuit sized for the case's 150 µm target
the gap is +0.2 points, and it runs from +10.0 to -20.8 points as the
assumed laboratory grind goes from 75 to 300 µm, which the open data do
not give. At each sample's own target-grind throughput it is +4.0
points, of which 3.8 come from the longer flotation residence at the
lower throughput rather than from the grind. At every assumed grind the
engine carries no information about which sample recovers more (Pearson
r between 0.03 and 0.06). It is a comparison, not a calibration: the
circuit, the breakage, the liberation and the flotation are authored for
another ore, and a locked-cycle test is not a plant.

\textbf{Plant hours.} One iron-ore plant's reverse flotation record
(Kaggle dataset 6294, CC0; 737,453 rows) gives hourly laboratory silica
beside 21 feed, reagent, pulp and column sensors. Dropping the 310 hours
whose silica label was interpolated leaves 3,701 pairs of consecutive
hours. Scored on three future windows after at least 24 hours of
embargo, the previous assay alone forecasts the next hour's silica with
a mean absolute error of 0.464 points, and the best sensor-only model,
ridge, is 0.001 points below the training mean's 0.766. In 461 of those
pairs the laboratory value is one carried over unchanged for three or
more consecutive hours; without them the previous assay's error is 0.510
points, still the lowest, and ridge is 0.014 points above the training
mean. The engine has no reverse cationic flotation family, so these
hours are shown and never simulated.

\section{Discussion and limitations}\label{discussion-and-limitations}

The test bounds transfer between authored plants. A thirteenth authored
plant is not a real one: the engine is a steady-state model with
authored parameters, no froth stability, no dynamics and no prices, and
its twelve plants share its structural assumptions. That is a feature of
the design, since the teacher's answer is exact and reproducible, and a
limit of the conclusion, since a real plant can depart from all twelve
at once.

Four findings should carry over to any surrogate of a flowsheet model.
First, interpolation and transfer are different questions, and the model
that answers the first best can answer the second worst. Second, holding
out one plant tests transfer only to the extent that no training plant
shares its ore: five copper plants on one mineral each look transferable
alone and transfer far worse when held out together, so the held-out
unit has to be the ore, not the plant. Third, transfer error is an
average over held-out plants that hides where it fails; the plants
unlike the others dominate it, and reporting the worst held-out case
matters as much as the mean. Fourth, a reconstruction guard misses
shifts in inputs that do not move with the others, so its false-accept
rate must be reported by input and by distance.

A calibrated study would replace the authored targets by measured
metallurgy across operating periods, calibrate the engine on some
periods and hold out an entire period and an ore family, and score both
point error and whether the surrogate preserves the ranking of operating
decisions.

\section{Reproducibility}\label{reproducibility}

The repository is https://github.com/fsantibanezleal/CAOS\_OreFlow
(MIT). The engine, the bake, the records, the browser port, the service
and the documentation are versioned together; the source, records and
computational environment are versioned with the software release.
\texttt{./scripts/setup.ps1} builds the environments and
\texttt{./scripts/precompute.ps1} regenerates every record. On a
workstation with 32 logical cores and an RTX 4070 Laptop GPU, the
GPU-accelerated final MLP fit is a computational component of a
simulator study over authored cases; it is not plant deployment or
evidence of plant-calibrated performance. \texttt{./scripts/smoke.ps1}
runs the checks: the guards, the Python and frontend suites and the
build, with the counts of each release in
\texttt{docs/release-verification.md}. The workbench at
https://oreflow.ml.fasl-work.com runs the engine in the browser on any
state of any case.

\section*{References}

\begin{itemize}
\sloppy
\raggedright
\tightlist
\item
  Audet, C. and Dennis, J.E. (2006). Mesh adaptive direct search
  algorithms for constrained optimization. SIAM Journal on Optimization
  17(1):188-217. doi:10.1137/040603371
\item
  Audet, C. and Dennis, J.E. (2009). A progressive barrier for
  derivative-free nonlinear programming. SIAM Journal on Optimization
  20(1):445-472. doi:10.1137/070692662
\item
  Bond, F.C. (1952). The third theory of comminution. Transactions AIME
  193:484-494.
\item
  Booker, A.J., Dennis, J.E., Frank, P.D., Serafini, D.B., Torczon, V.
  and Trosset, M.W. (1999). A rigorous framework for optimization of
  expensive functions by surrogates. Structural Optimization 17(1):1-13.
  doi:10.1007/BF01197708
\item
  Breiman, L. (2001). Random forests. Machine Learning 45:5-32.
  doi:10.1023/A:1010933404324
\item
  Duarte, R., Yamashita, A., da Silva, M., Cota, L. and Euzébio, T.
  (2021). Calibration and validation of a cone crusher model with
  industrial data. Minerals 11(11):1256. doi:10.3390/min11111256
\item
  Friedman, J.H. (2001). Greedy function approximation: a gradient
  boosting machine. Annals of Statistics 29(5):1189-1232.
  doi:10.1214/aos/1013203451
\item
  Global Mining Guidelines Group (2021). Determining the Bond Efficiency
  of Industrial Grinding Circuits, GMG01-MP-2021.
\item
  Gorain, B.K., Franzidis, J.-P. and Manlapig, E.V. (1997). Studies on
  impeller type, impeller speed and air flow rate in an industrial scale
  flotation cell. Part 4: Effect of bubble surface area flux on
  flotation performance. Minerals Engineering 10(4):367-379.
  doi:10.1016/S0892-6875(97)00014-9
\item
  Gorain, B.K., Franzidis, J.-P. and Manlapig, E.V. (1999). The
  empirical prediction of bubble surface area flux in mechanical
  flotation cells from cell design and operating data. Minerals
  Engineering 12(3):309-322. doi:10.1016/S0892-6875(99)00008-4
\item
  Herbst, J.A. and Fuerstenau, D.W. (1980). Scale-up procedure for
  continuous grinding mill design using population balance models.
  International Journal of Mineral Processing 7(1):1-31.
  doi:10.1016/0301-7516(80)90034-4
\item
  Herman, J. and Usher, W. (2017). SALib: an open-source Python library
  for sensitivity analysis. Journal of Open Source Software 2(9):97.
  doi:10.21105/joss.00097
\item
  Hoffimann, J., Augusto, J., Resende, L., Mathias, M., Mazzinghy, D.,
  Bianchetti, M. et al.~(2022). Modeling geospatial uncertainty of
  geometallurgical variables with Bayesian models and Hilbert-Kriging.
  Mathematical Geosciences 54(7):1227-1253.
  doi:10.1007/s11004-022-10013-1
\item
  King, R.P. (1979). A model for the quantitative estimation of mineral
  liberation by grinding. International Journal of Mineral Processing
  6:207-220. doi:10.1016/0301-7516(79)90037-1
\item
  Laplante, A.R. and Gray, S. (2005). Advances in gravity gold
  technology. Developments in Mineral Processing 15:280-307.
  doi:10.1016/S0167-4528(05)15013-3
\item
  Laplante, A.R. and Staunton, W.P. Gravity recovery of gold, an
  overview of recent developments (AMIRA P420B).
\item
  Loshchilov, I. and Hutter, F. (2019). Decoupled weight decay
  regularization. arXiv:1711.05101.
\item
  Moly-Cop Tools (Sepúlveda). BallSim\_Direct and BallParam\_Direct
  spreadsheets and documentation.
\item
  Muthaphuli, P. (2014). Production of pelletizing concentrates from
  Zandrivierspoort magnetite/haematite ore by magnetic separation.
  Journal of the Southern African Institute of Mining and Metallurgy
  114(7).
\item
  Oliveira, E.M. Quality prediction in a mining process: one iron-ore
  flotation plant, March to September 2017 (CC0). Kaggle dataset 6294.
  https://www.kaggle.com/datasets/edumagalhaes/quality-prediction-in-a-mining-process
\item
  Pereira, L., Frenzel, M., Khodadadzadeh, M., Tolosana-Delgado, R. and
  Gutzmer, J. (2021). A self-adaptive particle-tracking method for
  minerals processing. Journal of Cleaner Production 279:123711.
  doi:10.1016/j.jclepro.2020.123711
\item
  Plitt, L.R. (1976). A mathematical model of the hydrocyclone
  classifier. CIM Bulletin 69(776):114-123.
\item
  Polat, M. and Chander, S. (2000). First-order flotation kinetics
  models and methods for estimation of the true distribution of
  flotation rate constants. International Journal of Mineral Processing
  58:145-166. doi:10.1016/S0301-7516(99)00069-1
\item
  Rasmussen, C.E. and Williams, C.K.I. (2006). Gaussian Processes for
  Machine Learning. MIT Press.
\item
  Saltelli, A., Annoni, P., Azzini, I., Campolongo, F., Ratto, M. and
  Tarantola, S. (2010). Variance based sensitivity analysis of model
  output. Design and estimator for the total sensitivity index. Computer
  Physics Communications 181(2):259-270. doi:10.1016/j.cpc.2009.09.018
\item
  Savassi, O.N., Alexander, D.J., Franzidis, J.P. and Manlapig, E.V.
  (1998). An empirical model for entrainment in industrial flotation
  plants. Minerals Engineering 11(3):243-256.
  doi:10.1016/S0892-6875(98)00003-X
\item
  Torczon, V. (1997). On the convergence of pattern search algorithms.
  SIAM Journal on Optimization 7(1):1-25. doi:10.1137/S1052623493250780
\item
  Trahar, W.J. (1981). A rational interpretation of the role of particle
  size in flotation. International Journal of Mineral Processing
  8(4):289-327. doi:10.1016/0301-7516(81)90019-3
\item
  Vinnett, L. and Waters, K.E. (2025). The use of compressed
  exponentials for kinetic modelling of batch flotation. Minerals
  Engineering 226:109246. doi:10.1016/j.mineng.2025.109246
\end{itemize}
\end{document}
